\documentclass[11pt,a4paper]{article}
\usepackage[T1]{fontenc}
\usepackage[utf8]{inputenc}
\usepackage{lmodern}
\usepackage{amsmath,amssymb,amsthm,mathtools,mathrsfs}
\usepackage{graphicx,float,xcolor,multicol}
\usepackage[shortlabels]{enumitem}
\usepackage[margin=27mm]{geometry}
\usepackage{microtype,needspace,placeins}
\usepackage{tikz,tikz-cd}
\usetikzlibrary{arrows.meta,calc,positioning,patterns,decorations.pathreplacing}
\usepackage[colorlinks=true,linkcolor=teal,urlcolor=teal,citecolor=teal]{hyperref}
\setlength{\parindent}{0pt}
\setlength{\parskip}{.6em}
\setcounter{secnumdepth}{0}
\allowdisplaybreaks
\raggedbottom
\providecommand{\cross}{\times}
\DeclareUnicodeCharacter{2020}{\ensuremath{\dagger}}
\DeclareMathOperator{\supp}{supp}
\DeclareMathOperator*{\esssup}{ess\,sup}
\providecommand{\chapter}{\section}
\newenvironment{tool}[2]{\subsection*{Tool #1 --- #2}}{}
\newcommand{\ToolRef}[1]{Tool #1}
\newcommand{\FigureTag}[2]{\textit{Figure #1}}
\newcommand{\Result}[1]{\subsection*{#1}}
\newcommand{\Application}[1]{\subsubsection*{Application\if\relax\detokenize{#1}\relax\else\space--- #1\fi}}
\newcommand{\ProofHeading}{\par\textit{Proof.}\space}
\newcommand{\ProofOf}[1]{\par\textit{Proof of #1.}\space}
\newcommand{\RemarkHeading}{\par\textbf{Remark.}\space}
\newcommand{\NoteHeading}{\par\textbf{Note.}\space}
\newcommand{\ExerciseHeading}{\par\textbf{Exercise.}\space}

% Rendering support only; mathematical text is in each independent post.
\usepackage{booktabs,array,longtable}
\definecolor{HandoutBlue}{HTML}{2F6690}
\definecolor{HandoutNavy}{HTML}{17365D}
\newtheorem{theorem}{Theorem}
\newtheorem{lemma}[theorem]{Lemma}
\newtheorem{proposition}[theorem]{Proposition}
\newtheorem{corollary}[theorem]{Corollary}
\newtheorem{conjecture}[theorem]{Conjecture}
\newtheorem{definition}{Definition}
\newtheorem{example}{Example}
\newtheorem{namedtoolcounter}{Tool}
\newenvironment{namedtool}[1]{\begin{namedtoolcounter}[#1]}{\end{namedtoolcounter}}
\newtheorem*{claim}{Claim}
\newtheorem*{remark}{Remark}
\newtheorem*{exercise}{Exercise}
\newtheorem*{question}{Question}
\newtheorem*{observation}{Observation}
\newcommand{\SourceChapter}[2]{\section*{#2}}
\newcommand{\Lecture}[2]{\section*{Lecture #1: #2}}
\newcommand{\unclear}[1]{\textit{[unclear: #1]}}
\newcommand{\sourcegap}{\textit{[unfinished in the source]}}
\DeclareMathOperator{\dist}{dist}
\DeclareMathOperator{\id}{id}
\DeclareMathOperator{\Hol}{Hol}
\DeclareMathOperator{\Per}{Per}
\DeclareMathOperator{\Fix}{Fix}
\DeclareMathOperator{\diam}{diam}
\DeclareMathOperator{\Var}{Var}
\DeclareMathOperator{\Leb}{Leb}
\DeclareMathOperator{\Hom}{Hom}
\DeclareMathOperator{\End}{End}
\DeclareMathOperator{\Ann}{Ann}
\DeclareMathOperator{\Spec}{Spec}
\DeclareMathOperator{\Max}{Max}
\DeclareMathOperator{\Ob}{Ob}
\DeclareMathOperator{\Tor}{Tor}
\DeclareMathOperator{\Ass}{Ass}
\DeclareMathOperator{\Mor}{Mor}
\DeclareMathOperator{\Fun}{Fun}
\DeclareMathOperator{\Nat}{Nat}
\DeclareMathOperator{\Cone}{Cone}
\DeclareMathOperator{\MCG}{MCG}
\DeclareMathOperator{\Aut}{Aut}
\DeclareMathOperator{\Deck}{Deck}
\DeclareMathOperator{\SL}{SL}
\DeclareMathOperator{\PSL}{PSL}
\DeclareMathOperator{\Area}{Area}
\DeclareMathOperator{\im}{im}
\DeclareMathOperator{\PSh}{PSh}
\newcommand{\CRing}{\mathsf{CRings}}
\newcommand{\AMod}{A\text{-}\mathsf{Mod}}
\newcommand{\Ab}{\mathsf{Ab}}
\newcommand{\Z}{\mathbb Z}
\newcommand{\Q}{\mathbb Q}
\newcommand{\R}{\mathbb R}
\newcommand{\C}{\mathbb C}
\newcommand{\N}{\mathbb N}
\newcommand{\kk}{\mathbb k}
\renewcommand{\aa}{\mathfrak a}
\newcommand{\bb}{\mathfrak b}
\newcommand{\pp}{\mathfrak p}
\newcommand{\qq}{\mathfrak q}
\newcommand{\mm}{\mathfrak m}
\newcommand{\nn}{\mathfrak n}
\newcommand{\rad}{\mathrm r}
\newcommand{\cat}[1]{\mathsf{#1}}
\setcounter{secnumdepth}{0}
\allowdisplaybreaks

\graphicspath{{figures/commutative-algebra/}}
\begin{document}
\section*{The Nullstellensatz and Affine Geometry}
% BLOG-CONTENT-BEGIN
\subsection{Hilbert's Nullstellensatz}
% Source page 5
\setcounter{definition}{14}\begin{definition}
A ring $R$ is \emph{Jacobson} if, for every prime $\pp$,
\[\pp=\bigcap_{\substack{\mm\supseteq\pp\\\mm\text{ maximal}}}\mm.\]
\end{definition}
\setcounter{example}{8}\begin{example}
\begin{enumerate}[label=(\roman*)]
\setcounter{enumi}{0}\item A field $\kk$ is a Jacobson ring.
\setcounter{enumi}{1}\item $\Z$ is Jacobson: every nonzero prime ideal is maximal, and $(0)=\operatorname{nilradical}(\Z)=\bigcap_{\mm\text{ maximal}}\mm=\bigcap_{\pp\text{ prime}}\pp$.
\end{enumerate}
\end{example}
The following lemma characterizes Jacobson rings.
\setcounter{theorem}{23}\begin{lemma}
$R$ is Jacobson if and only if, for prime $\pp$ and $0\ne b\in R/\pp$, whenever $(R/\pp)_{(b)}$ is a field, $R/\pp$ is a field.
\end{lemma}
\setcounter{theorem}{24}\begin{theorem}
Let $R$ be a Jacobson ring and $S$ a finitely generated $R$-algebra. Then:
\begin{enumerate}[label=(\roman*)]
\setcounter{enumi}{0}\item $S$ is a Jacobson ring.
\setcounter{enumi}{1}\item If $\nn<S$ is maximal, then $\mm:=\nn\cap R$ is maximal in $R$, and $S/\nn$ is a finite field extension of $R/\mm$.
\end{enumerate}
\end{theorem}
\setcounter{theorem}{25}\begin{lemma}
Let $A\subseteq B$ be domains, with $B$ integral over $A$. Then $B$ is a field if and only if $A$ is a field.
\end{lemma}
\setcounter{theorem}{27}\begin{corollary}
Let $\kk$ be algebraically closed. Every maximal ideal of $\kk[x_1,\ldots,x_n]$ has the form $\mm_P=(x_1-a_1,\ldots,x_n-a_n)$ for $P=(a_1,\ldots,a_n)\in\kk^n$.
\end{corollary}
\begin{proof}
Since $\kk[x_1,\ldots,x_n]/\mm_P\cong\kk$, $\mm_P$ is maximal. If $\nn$ is maximal in $\kk[x_1,\ldots,x_n]$, then $\kk[x_1,\ldots,x_n]/\nn$ is algebraic over $\kk/(\nn\cap\kk)=\kk$. Since $\kk$ is algebraically closed, $\kk[x_1,\ldots,x_n]/\nn=\kk$. Let $a_i=[x_i]$ in this quotient and $P=(a_1,\ldots,a_n)$. Then $\mm_P\subseteq\nn$. Since $\mm_P$ is maximal, $\nn=\mm_P$.
\end{proof}
\subsection{A little geometry}
Consider the two objects $\kk^n$ and $\kk[x_1,\ldots,x_n]$. To $I\subseteq\kk[x_1,\ldots,x_n]$ associate
\[Z(I):=\{P\in\kk^n\mid f(P)=0\text{ for every }f\in I\}.\]
% Source page 10
Clearly $Z(I)=Z((I))$. Sets of this form are \emph{algebraic sets} of $\kk^n$. For ideals $\{J_i\}_i$ of $\kk[x_1,\ldots,x_n]$,
\[
\bigcap_i Z(J_i)=Z\left(\bigcup_i J_i\right),\qquad
\bigcup_{i=1}^n Z(J_i)=Z\left(\prod_{i=1}^n J_i\right).
\]
On an algebraic set $X$, declaring algebraic subsets of $X$ to be closed defines the \emph{Zariski topology}.

For a subset $X\subseteq\kk^n$, set $I(X):=\{f\in\kk[x_1,\ldots,x_n]\mid f(P)=0\text{ for every }P\in X\}$. Clearly $I(X)$ is an ideal. For algebraic $X$, define its \emph{coordinate ring}, also called an \emph{affine ring}, by $A(X):=\kk[x_1,\ldots,x_n]/I(X)$.

If $f\in A(X)$ and $f^n\equiv0$, then $f^n(P)=0$ for all $P\in X$. Since $\kk$ is a domain, $f\equiv0$. Thus $\{0\}$ is the only nilpotent element in $A(X)$, and $I(X)=\rad(I(X))$. Such ideals are called \emph{radical ideals}.
\begin{remark}
Clearly $Z(I(X))=X$. Thus we have the correspondences
\[
\begin{aligned}
\text{algebraic sets}&\longleftrightarrow\text{radical ideals}\\
&\longleftrightarrow\text{quotients of }\kk[x_1,\ldots,x_n]\text{ by radical ideals}\\
&\longleftrightarrow\text{coordinate rings}.
\end{aligned}
\]
\end{remark}
\setcounter{definition}{15}\begin{definition}
Let $X\subseteq\kk^m$ and $Y\subseteq\kk^n$ be algebraic sets. A map
\[
F:\kk^m\to\kk^n,\qquad
(a_1,\ldots,a_m)\mapsto(f_1(a_1,\ldots,a_m),\ldots,f_n(a_1,\ldots,a_m)),
\]
where $f_i\in\kk[x_1,\ldots,x_m]$, is a \emph{morphism of $X$ to $Y$} if $F|_X:X\to Y$.
\end{definition}
Given $F:X\to Y$, we have
\[
F^{\#}:\kk[y_1,\ldots,y_n]\longrightarrow\kk[x_1,\ldots,x_m],
\qquad y_i\longmapsto f_i(x_1,\ldots,x_m).
\]
% BLOG-CONTENT-END
\end{document}
